\documentclass[11pt,a4paper]{article}
\usepackage[T1]{fontenc}
\usepackage{amsmath,amssymb}
\usepackage[margin=1in]{geometry}
\usepackage{hyperref}
\usepackage{xurl}
\title{Progressive-Lens-Inspired Distortion in VR}
\author{Quest Shooting Demo}
\date{}
\begin{document}
\maketitle

\section{Scope and safety}
This project applies a controllable, \emph{illustrative} image warp to the
rendered virtual scene. It is inspired by the approach in Stein et al.,
``Simulating Lens Distortion in Virtual Reality'' (2020), not a reproduction
of their measured progressive addition lens (PAL). The paper's fitted
coefficients and measured displacement map are not supplied in the paper
images. The constants in this implementation are design choices, not
calibration data, clinical measurements, or verified Quest optics. The
simulation starts off. Stop using it if you experience discomfort.

\section{What: optics and the paper's method}
A progressive addition lens varies optical power across its surface, providing
regions intended for distance and near viewing with a transition corridor.
Outside that corridor the lens can introduce spatially varying optical
aberrations and image displacement. The mapping from the scene to the retinal
image can therefore vary across the visual field; it is not generally a
simple, rotationally symmetric camera-lens distortion.

Stein et al.\ measured a lens distortion map on a grid, fit two polynomial
functions for horizontal and vertical displacement, and used a pixel shader to
sample the undistorted scene at displaced positions. Their experiment applied
the same distortion to both eyes. They reported a statistically significant
effect on heading estimates in their study; that finding does not validate our
illustrative model. The paper describes a fifth-order polynomial family,
whereas the formulas below are a deliberately simple third-order example and
are not the authors' fitted functions.

\section{How: the implemented displacement}
For each eye's screen texture, let $\mathbf{q}=(u,v)$ be the output pixel's
normalized texture coordinates, with $(0,0)$ at the upper-left and $(1,1)$ at
the lower-right. Define
\begin{equation}
 x=2u-1,\qquad y=2v-1,\qquad L=\frac{y+1}{2}.
\end{equation}
Here $x$ and $y$ run from $-1$ to $1$, and $L$ grows from zero at the top
to one at the bottom. The implemented source-coordinate displacements are
\begin{align}
 \Delta u &= s(0.08)xL^2,\\
 \Delta v &= s(0.04)x^2(1+y),
\end{align}
where $s$ is the user-controlled dimensionless strength. The source position
and output color are
\begin{equation}
 \mathbf{q}_{\mathrm{src}}=
 \operatorname{clamp}\bigl(\mathbf{q}+(\Delta u,\Delta v),\,0.001,\,0.999\bigr),
 \qquad C_{\mathrm{out}}(\mathbf{q})=C_{\mathrm{in}}(\mathbf{q}_{\mathrm{src}}).
\end{equation}
The clamp is component-wise and avoids reading beyond the image; near an edge
it can stretch the border. This is a \emph{backward} or pull sampling map:
the displacement identifies where an output pixel reads from, not where a
source feature moves. Visible feature motion is therefore generally opposite
the sampling displacement. The same UV-space formula is used for both eyes; it
does not use gaze tracking or independently measured per-eye lens maps.

For example, at $(u,v)=(0.75,0.75)$, $(x,y)=(0.5,0.5)$ and $L=0.75$.
The displacement is $s(0.0225,0.015)$. At $s=0.2$, the source coordinate is
shifted by $(0.0045,0.003)$ in texture units.
The effect is off on launch. Right-controller A hides or shows the panel and B
interacts with it. The on/off button controls the effect; strength runs from
0.0 to 1.0 in steps of 0.1, and reset restores 0.2. A zero strength is an
identity mapping except at the clamped image boundary.

\section{Why: optical interpretation and limitations}
The horizontal function is odd in $x$, so its sampling displacement has
opposite signs on the left and right sides. The vertical function is even in
$x$, producing matching vertical sampling displacement on both sides. Both
grow toward the lower image and vanish on the centerline, giving a
qualitatively asymmetric, peripheral warp inspired by a progressive lens's
nonuniform field. The factors 0.08 and 0.04 set the relative magnitudes of
these illustrative terms; $s$ scales both together. They have not been
inferred from a physical lens.

The shader does not trace rays through an optical prescription, simulate
depth-dependent refraction, blur or astigmatic point-spread functions, eye
rotation, or a changing gaze position. It modifies the opaque rendered scene
texture, not the headset's real-world passthrough feed; transparent objects
and UI may not be included. A PAL-specific map requires measured or
ray-traced data, the lens/eye geometry, and careful coordinate conventions.
An accurate stereo reproduction can require different maps for each eye.
Neither fitting a polynomial nor adding a screen warp establishes perceptual
or clinical validity; validate against independent measurements, both eyes,
and user comfort.

\section{Reference}
N. Stein, K. Rifai, S. Wahl, and M. Lappe, ``Simulating Lens Distortion in
Virtual Reality,'' Proceedings of the 13th ICDVRAT with ITAG, 2020,
\url{https://niklas-hypki.de/Files/LensDistortion.pdf}.
\end{document}
